\section{Discussion}\label{sec:discussion}

\subsection{Why other studies reach different conclusions}
The results differ from widely cited cost comparisons for identifiable reasons. \citet{Lazard2025} reports 175--255\,USD/MWh for US nuclear, derived from Vogtle 3 and 4, the first US units in thirty years, at 12{,}300--17{,}570\,USD/kW and 7.7\% cost of capital: a first-of-a-kind American cost. \citet{Fraunhofer2024} reports 136--490\,EUR/MWh for Germany assuming 4{,}300--6{,}300 operating hours today and 2{,}000--4{,}000 by 2045, i.e.\ a reactor idle for more than half of the year, together with 6{,}000--16{,}000\,EUR/kW and 7.8\% real cost of capital; in our central PI system nuclear runs about <<UN:ss0_PI_670_h3:NUC:hours|0>> equivalent hours (<<EXPR:V('UN:ss0_PI_670_h3:NUC:hours')/87.6|0>>\%), and the stand-alone cost at 90\% (Table~\ref{tab:lcoe}) understates its cost per MWh slightly. The same report gives 6--22.5\,ct/kWh for solar with batteries, for systems that shift part of their own output by a few hours without any guarantee of covering load. \citet{Lund2026} find renewables 53\% cheaper than nuclear in Denmark with an effective nuclear cost of 10{,}000\,EUR/kW and 8\% cost of capital. The 188--991\,USD/MWh reported by Banca d'Italia \citep{Bankitalia2025} comes from a study of small modular reactors \citep{Steigerwald2023}, not of the large reactors considered here, and the Italia Solare scenarios assume 11{,}000\,EUR/kW in 2040 \citep{ItaliaSolare2026}, above the Dukovany contract. Their Vogtle-based figures are our risk cases, not the central case.

System studies without nuclear rely on technologies without market prices for the winter. \citet{Lombardi2020} cover low-wind, low-sun periods with existing gas plants burning synthetic methane from electrolysis and direct air capture, which would require about 3.5 billion m$^3$ per year (our estimate) against 0.4--0.57 billion m$^3$ of biomethane produced today \citep{Biometano2024}. \citet{Franzmann2026} obtain 91.6\,EUR/MWh for a fully renewable Europe with 1{,}116\,GW of electrolysers and 501\,GW of hydrogen turbines, hydrogen accounting for 22\% of system cost, without sensitivity analysis on hydrogen or grid costs. The study closest to ours, \citet{Mastrantuono2026}, finds nuclear competitive, but its 8\,GWe cap binds; our model, without a cap, at that cost (about 3{,}600\,EUR/kW at 1.1188\,USD per euro, without inflation, with 8 years of construction: fixN 373 at 5\% and 488 at 7\%, our calculation) chooses about <<N6:PI:s:5:373|0>>\,GW at 5\% (PI, six hours, interpolated) and at 7\% more than the <<RUN:ss0_PI_7pc_700_h6:N_GW|0>>\,GW of the first computed point (fixN 700), far above 8\,GWe, so the cap would also bind in our model. The unpublished Politecnico di Milano analysis \citep{PoliMi2026} finds up to 15\,GW of nuclear at 5{,}500\,EUR/kW and none at 8{,}300\,EUR/kW with ten-year construction. We do not know its discount rate. With seven years, 5{,}500\,EUR/kW is fixN 485 at 5\% and 649 at 7\%, and our model gives <<N6:PI:s:5:485|0>>\,GW (PI, six hours) at 5\% and more than <<RUN:ss0_PI_7pc_700_h6:N_GW|0>> at 7\%, above their 15\,GW; with ten years, 8{,}300\,EUR/kW is fixN 696 at 5\% and 989 at 7\%, and our model gives <<N6:PI:s:5:696|1>>\,GW (PI) and <<N6:OB:s:5:696|1>> (OB) at 5\% and <<N6:PI:s:7:989|1>> and <<N6:OB:s:7:989|1>> at 7\%, so under net self-sufficiency their conclusion of no nuclear does not coincide with ours in any case. In the central PI case at three hours a system without nuclear costs <<EXPR:(V('RUN:ss0_PI_none_h3:cost_new')/V('RUN:ss0_PI_670_h3:cost_new')-1)*100|0>>\% more than the system with nuclear (<<RUN:ss0_PI_none_h3:cost_new|2>> against <<RUN:ss0_PI_670_h3:cost_new|2>>\,EUR/MWh) and in OB <<EXPR:(V('RUN:ss0_OB_none_h3:cost_new')/V('RUN:ss0_OB_670_h3:cost_new')-1)*100|0>>\% more (<<RUN:ss0_OB_none_h3:cost_new|2>> against <<RUN:ss0_OB_670_h3:cost_new|2>>): both are below the 11--163\% increase implied by \citet{Sepulveda2018} (our conversion of their 10--62\% reduction), and the gaps in EUR/MWh (<<GAP:ss0_PI_none_h3:ss0_PI_670_h3|2>> and <<GAP:ss0_OB_none_h3:ss0_OB_670_h3|2>>) are in the low range of the 8--50\,USD/MWh integration costs of \citet{NEA2019}, because the model has imports, hydrogen and other firm sources. Without imports at five per cent and six hours, with nuclear at fixN 670, the increase is <<EXPR:(120.17/V('C6:PI:n:5:670')-1)*100|0>>\% (PI, 120.17 against <<C6:PI:n:5:670|2>>) and <<EXPR:(101.57/V('C6:OB:n:5:670')-1)*100|0>>\% (OB, 101.57 against <<C6:OB:n:5:670|2>>), with costs interpolated between computed points. The comparison is approximate because systems and costs differ.

\citet{Pastore2025} describe a fully renewable Italy in 2050 with an hourly model, assuming a potential of 115.4\,GW of onshore and 55.7\,GW of offshore wind, and install the whole potential in most scenarios. Our central onshore limit is 69.6\,GW, in blocks of 1{,}200--2{,}400 full-load hours at rising cost; the JRC estimate under current siting rules is 96.9\,GW \citep{JRCEUTIMES}. Their result therefore rests on sites beyond our curve. In our model the central PI system builds <<GR:ss0_PI_670_h6:on:GW|1>>\,GW of onshore wind at six hours (<<GR:ss0_PI_670_h3:on:GW|1>> at three) and leaves about 40\,GW of blocks E3-E5 unused, at costs rising to 118\,EUR/MWh: a longer curve at rising cost would lower mainly the cost of the system without nuclear.

Models that impose net-zero rather than zero direct emissions can cover winter with gas plants fitted with carbon capture and storage and offset the residual with direct air capture. No gas-fired power plant with capture is yet in operation: the first, Net Zero Teesside in the UK (742\,MW, up to about 2\,Mt CO$_2$ a year), is expected in 2028, and its contract terms are not public \citep{NZTPower}. The only power plant with capture and years of operation, the coal-fired Boundary Dam 3 in Canada, captured on average 57\% of its CO$_2$ against a 90\% target, with capture running in 80\% of operating hours \citep{IEEFABD3}. In Italy, Ravenna CCS targets 4\,Mt a year by 2030 from industrial emitters \citep{EniRavenna2026}. The most favourable full-chain cost estimate, about GBP\,110 per tonne, adds about GBP\,37 per MWh, and other estimates are twice as high \citep{Goodall2025}. At early-October 2026 gas prices (PSV about 74\,EUR/MWh, \citealp{SelectraPSV2026}) the fuel alone of a combined cycle with capture costs about 150\,EUR per MWh of electricity (our calculation at 50\% efficiency), against <<RUN:ss0_PI_670_h3:cost_new|1>>\,EUR/MWh for the whole system with nuclear in the central case at three hours. In the model gas exists only in the NZ variants and costs 38.5\,EUR/MWh, the 2025 average PSV \citep{GME2025}, about half the autumn 2026 price, so the variant favours gas. Direct air capture is further away: the largest plant, Mammoth in Iceland, captured 675\,t in the first half of 2026 against a nameplate of 36{,}000\,t a year \citep{Climeworks2026}. A low-nuclear result that depends on these technologies holds only if they reach the assumed cost and performance. This can be checked after 2028: if Teesside shows at least 90\% capture at costs near GBP\,100 per tonne, the gap of nuclear should be recomputed with the NZ variants of Section~\ref{sec:sens}.

\subsection{Construction cost and time}
The nuclear share depends on the annualised fixed cost, hence on three variables that policy can influence. Construction cost falls with series production of one design: domestic Korean units cost 2{,}400--3{,}300\,USD/kW, the exported Barakah units 3{,}600--4{,}400\,USD/kW, the Dukovany contract, 8{,}540\,EUR/kW in the model (EPC price net of fuel, plus owner's costs) \citep{HSF2026,KHNP2025,WNAuae,Dukovany2024}. Cost of capital falls with public financing: the European Commission approved for Dukovany a subsidised state loan covering most construction costs and a 40-year contract for difference \citep{WNN2024}. Construction time falls with completed engineering and stable rules. Flamanville, at about 10{,}300\,EUR/kW in 2025 euros \citep{EurostatHICP} (8{,}000 in 2015 euros), took about 17 years \citep{Flamanville2024} and Hinkley Point C will take about 13 \citep{Hinkley2026}, but the median for reactors started since 1990 is 5.7 years and 89\% were completed in less than 10 \citep{Ritchie2023}. For Flamanville, the Folz report points to industrial causes, incomplete engineering and loss of skills \citep{Folz2019}; in the United States, rules changed during construction after Three Mile Island raised costs and durations \citep{Lovering2016}.


\subsection{Deployment rate and siting}
The central optimum, <<RUN:ss0_PI_670_h3:N_GW|1>>\,GW at three hours (PI), corresponds to about <<EXPR:V('RUN:ss0_PI_670_h3:N_GW')/1.055|0>> APR1000 units of 1{,}055\,MWe, the unit of the central case, or about <<EXPR:V('RUN:ss0_PI_670_h3:N_GW')/1.4|0>> APR1400 of 1.4\,GW. This number is a constraint to verify, not a fact: the model describes a regime system with 2040 demand and all fleets already built, not the situation in 2040. With a first unit in 2035, <<EXPR:V('RUN:ss0_PI_670_h3:N_GW')/1.055|0>> APR1000 units by 2050 require about <<EXPR:V('RUN:ss0_PI_670_h3:N_GW')/1.055/15|1>> a year for 15 years, about <<EXPR:V('RUN:ss0_PI_670_h3:N_GW')/15|1>>\,GW a year (<<RUN:ss0_PI_670_h3:N_GW|1>> over 15). According to the IAEA reactor tables, France connected 52 reactors to the grid between 1977 and 1990, 54.4\,GW net: 3.9\,GW a year on average, with peaks of 8 reactors in 1981 and 7.4\,GW in 1986 \citep{IAEAFrance}. Italy would need about <<EXPR:V('RUN:ss0_PI_670_h3:N_GW')/15/3.9*100|0>>\% of the average French pace of that period (<<EXPR:V('RUN:ss0_PI_670_h3:N_GW')/15|1>> over 3.9). Korea commissioned about one reactor a year in the 1980s \citep{WNAcountries}. The required pace is below the French and similar to the Korean one, in a country restarting without a full supply chain. The model does not compute the cost of a slower pace: if the pace stays below <<EXPR:V('RUN:ss0_PI_670_h3:N_GW')/1.055/15|1>> reactors a year the steady-state fleet is reached later, and wind, solar, imports and batteries can be added meanwhile.

Siting differs sharply between the options. Reactors are usually built in groups of four to six per site, so <<EXPR:V('RUN:ss0_PI_670_h3:N_GW')/1.055|0>> APR1000 units require <<EXPR:V('RUN:ss0_PI_670_h3:N_GW')/1.055/6|0>> to <<EXPR:V('RUN:ss0_PI_670_h3:N_GW')/1.055/4|0>> sites, which could include former thermal plants already connected to the extra-high-voltage grid if they meet seismic and cooling requirements. The central system (PI, six hours) needs <<EXPR:V('GR:ss0_PI_670_h6:on:GW')-13.6|1>>\,GW of new onshore wind, about <<EXPR:(V('GR:ss0_PI_670_h6:on:GW')-13.6)*200|-1>> turbines of 5\,MW in <<EXPR:(V('GR:ss0_PI_670_h6:on:GW')-13.6)*10|-1>>--<<EXPR:(V('GR:ss0_PI_670_h6:on:GW')-13.6)*20|-1>> parks, and <<EXPR:V('UN:ss0_PI_670_h6:PV_terra_Sud-isole:GW')+V('UN:ss0_PI_670_h6:PV_terra-pipeline_Lazio/Sicilia/Piem./Puglia/Sard.:GW')+V('UN:ss0_PI_670_h6:PV_terra_Centro-Sud:GW')+V('UN:ss0_PI_670_h6:PV_terra_Nord:GW')|1>>\,GW of new ground-mounted solar, about <<EXPR:(V('UN:ss0_PI_670_h6:PV_terra_Sud-isole:GW')+V('UN:ss0_PI_670_h6:PV_terra-pipeline_Lazio/Sicilia/Piem./Puglia/Sard.:GW')+V('UN:ss0_PI_670_h6:PV_terra_Centro-Sud:GW')+V('UN:ss0_PI_670_h6:PV_terra_Nord:GW'))*1700|-2>> hectares at 1{,}700\,ha/GW \citep{ISPRA2026}, plus <<UN:ss0_PI_670_h6:PV_agrivoltaico_Sud-isole:GW|1>>\,GW of agrivoltaic, about <<EXPR:V('UN:ss0_PI_670_h6:PV_agrivoltaico_Sud-isole:GW')*2000|-2>> hectares. The system without nuclear requires <<EXPR:V('GR:ss0_PI_none_h6:on:GW')-13.6|1>>\,GW of new onshore wind, about <<EXPR:(V('GR:ss0_PI_none_h6:on:GW')-13.6)*200|-2>> turbines in <<EXPR:(V('GR:ss0_PI_none_h6:on:GW')-13.6)*10|-1>>--<<EXPR:(V('GR:ss0_PI_none_h6:on:GW')-13.6)*20|-1>> parks, and <<EXPR:V('GR:ss0_PI_none_h6:pv:GW')-24.2|1>>\,GW of new solar: <<EXPR:(V('UN:ss0_PI_none_h6:PV_terra_Sud-isole:GW')+V('UN:ss0_PI_none_h6:PV_terra-pipeline_Lazio/Sicilia/Piem./Puglia/Sard.:GW')+V('UN:ss0_PI_none_h6:PV_terra_Centro-Sud:GW')+V('UN:ss0_PI_none_h6:PV_terra_Nord:GW'))|1>>\,GW ground-mounted (about <<EXPR:(V('UN:ss0_PI_none_h6:PV_terra_Sud-isole:GW')+V('UN:ss0_PI_none_h6:PV_terra-pipeline_Lazio/Sicilia/Piem./Puglia/Sard.:GW')+V('UN:ss0_PI_none_h6:PV_terra_Centro-Sud:GW')+V('UN:ss0_PI_none_h6:PV_terra_Nord:GW'))*1700|-2>> hectares), <<EXPR:(V('UN:ss0_PI_none_h6:PV_agrivoltaico_Sud-isole:GW')+V('UN:ss0_PI_none_h6:PV_agrivoltaico_Centro-Sud:GW')+V('UN:ss0_PI_none_h6:PV_agrivoltaico_Nord:GW'))|1>>\,GW of agrivoltaic (about <<EXPR:(V('UN:ss0_PI_none_h6:PV_agrivoltaico_Sud-isole:GW')+V('UN:ss0_PI_none_h6:PV_agrivoltaico_Centro-Sud:GW')+V('UN:ss0_PI_none_h6:PV_agrivoltaico_Nord:GW'))*2000|-3>> hectares at 0.5\,MW per hectare) and <<EXPR:V('GR:ss0_PI_none_h6:pv:GW')-24.2-(V('UN:ss0_PI_none_h6:PV_terra_Sud-isole:GW')+V('UN:ss0_PI_none_h6:PV_terra-pipeline_Lazio/Sicilia/Piem./Puglia/Sard.:GW')+V('UN:ss0_PI_none_h6:PV_terra_Centro-Sud:GW')+V('UN:ss0_PI_none_h6:PV_terra_Nord:GW'))-(V('UN:ss0_PI_none_h6:PV_agrivoltaico_Sud-isole:GW')+V('UN:ss0_PI_none_h6:PV_agrivoltaico_Centro-Sud:GW')+V('UN:ss0_PI_none_h6:PV_agrivoltaico_Nord:GW'))|1>>\,GW of industrial rooftops; ground and agrivoltaic together are about <<EXPR:((V('UN:ss0_PI_none_h6:PV_terra_Sud-isole:GW')+V('UN:ss0_PI_none_h6:PV_terra-pipeline_Lazio/Sicilia/Piem./Puglia/Sard.:GW')+V('UN:ss0_PI_none_h6:PV_terra_Centro-Sud:GW')+V('UN:ss0_PI_none_h6:PV_terra_Nord:GW'))*1700+(V('UN:ss0_PI_none_h6:PV_agrivoltaico_Sud-isole:GW')+V('UN:ss0_PI_none_h6:PV_agrivoltaico_Centro-Sud:GW')+V('UN:ss0_PI_none_h6:PV_agrivoltaico_Nord:GW'))*2000)/1000|-1>> thousand hectares, about <<EXPR:((V('UN:ss0_PI_none_h6:PV_terra_Sud-isole:GW')+V('UN:ss0_PI_none_h6:PV_terra-pipeline_Lazio/Sicilia/Piem./Puglia/Sard.:GW')+V('UN:ss0_PI_none_h6:PV_terra_Centro-Sud:GW')+V('UN:ss0_PI_none_h6:PV_terra_Nord:GW'))*1700+(V('UN:ss0_PI_none_h6:PV_agrivoltaico_Sud-isole:GW')+V('UN:ss0_PI_none_h6:PV_agrivoltaico_Centro-Sud:GW')+V('UN:ss0_PI_none_h6:PV_agrivoltaico_Nord:GW'))*2000)/30.1e6*100|1>>\% of the 30.1 million hectares of the national territory (our estimates; with agrivoltaics farming continues). Individual wind projects meet strong local opposition: in Sardinia a regional moratorium and a law on suitable areas, later struck down by the Constitutional Court \citep{Sardegna2024,Consulta2026}; in the central Apennines 21 projects with 188 turbines contested by local committees \citep{Appennino2025}. Nuclear sites are fewer but each is more contested: in a 2026 Ipsos survey 12\% of Italians would accept a plant within 10\,km and 34\% within 100\,km \citep{Ipsos2026}. A few large consents are administratively simpler than thousands of small ones, but not politically easy. The national repository for radioactive waste, about 95{,}000\,m$^3$ including waste from medicine, research and industry, is needed in any case. It would occupy about 150 hectares, about as much as a single 80\,MW ground-mounted solar park \citep{Deposito2023}: one site, against <<EXPR:(V('GR:ss0_PI_670_h6:on:GW')-13.6)*10|-1>>--<<EXPR:(V('GR:ss0_PI_670_h6:on:GW')-13.6)*20|-1>> wind parks in the central system and <<EXPR:(V('GR:ss0_PI_none_h6:on:GW')-13.6)*10|-1>>--<<EXPR:(V('GR:ss0_PI_none_h6:on:GW')-13.6)*20|-1>> without nuclear. 51 potentially suitable areas have been identified \citep{Deposito2023}, and the only municipality that volunteered withdrew in 2024 \citep{Trino2024}.

\subsection{Supply security and materials}
Fuel is a small and storable share of nuclear cost. A light-water reactor uses about 0.021--0.027\,kg of natural uranium per MWh \citep{WNAecon}; at the long-term price of 96.5\,USD/lb U$_3$O$_8$ \citep{Cameco2026} this is about 5--7\,USD/MWh, and the whole fuel cycle about 11\,USD/MWh, roughly 5\% and 10--12\% of the cost of a new nuclear MWh. Doubling the uranium price raises the cost of a nuclear MWh by about 5\%, whereas fuel dominates the cost of gas-fired power. Production is spread over several countries (Kazakhstan 39\%, Canada 24\%, Namibia 12\%, Australia 8\% in 2024), although Kazakhstan, Uzbekistan and Russia together account for about half \citep{WNAmining}. The critical bottleneck is enrichment, where Rosatom holds 43\% of world capacity and Russia supplied 23\% of the enrichment services delivered to the EU in 2025 \citep{WNAenrich,ESA2025}. Italy has small uranium deposits: according to historical estimates, about 6{,}000\,t of uranium oxide (U$_3$O$_8$) at Val Vedello and about 1{,}000\,t at Novazza \citep{Geopop2026}, together about 5{,}900\,t of uranium, enough for roughly 30\,GW-years at 200\,tU per GW-year. Uranium from seawater is a long-run ceiling on the fuel price: published estimates of 606--1{,}230\,USD/kg \citep{SchneiderGill2014,Schneider2013}, 2.4--5 times today's market price, would add 7--26\,USD/MWh, up to about 23\,EUR/MWh (26 over 1.1188, our calculation), more than the gap of the central case (<<GAP:ss0_PI_none_h3:ss0_PI_670_h3|2>> in PI, <<GAP:ss0_OB_none_h3:ss0_OB_670_h3|2>> in OB, three hours); it is a theoretical ceiling on the price of uranium, not a forecast; the technology has been demonstrated at pilot scale but not commercially. One year of fuel for 1\,GW is about 200\,t of uranium, which is easily stockpiled; EU utilities already hold about 3.6 years of requirements \citep{ESA2025}.

Solar and batteries carry a different dependence. China holds more than 80\% of every stage of the solar supply chain and about 95\% of polysilicon and wafers \citep{IEApv2022}; 98\% of the solar panels imported into the EU in 2024 came from China \citep{Eurostat2025}. China refines about 70\% of lithium and more than 95\% of battery-grade graphite, and produces about 80\% of battery cells \citep{IEAgcmo2025}; it has imposed export controls on graphite, rare earths and battery technologies since 2023, partly suspended until November 2026 \citep{HSFchina2025}. Per MW, onshore wind requires about 10{,}100\,kg of critical minerals, solar 6{,}800\,kg and nuclear 5{,}300\,kg \citep{WNAminerals}; per MWh, given the load factors, nuclear requires several times less. These dependences matter more for the system without nuclear, which in the PI case at six hours needs <<GR:ss0_PI_none_h6:battery:GWh|0>>\,GWh of batteries and <<GR:ss0_PI_none_h6:pv:GW|0>>\,GW of solar, than for the central case, with <<GR:ss0_PI_670_h6:battery:GWh|0>>\,GWh and <<GR:ss0_PI_670_h6:pv:GW|0>>\,GW; the cost of reducing these dependences, for example with components meeting origin requirements, is not quantified by the present model. Nuclear has its own concentrations, in enrichment, fuel fabrication and heavy forgings, which only a handful of firms produce.


\subsection{Industrial base}
Large reactors are typically supplied by a foreign vendor responsible for the design and system integration, with local firms supplying components and balance of plant, as for Dukovany and Barakah. Italy has a supply chain for this role: more than 70 firms with about 2{,}800 direct jobs in nuclear activities in 2022 \citep{TEHA2024}, the refurbishment of Cernavoda unit 1 in Romania led by Ansaldo Nucleare, worth more than 400 million euros \citep{Cernavoda2024}, and the ITER vacuum vessel sectors, the last of which was delivered in October 2026 \citep{ITER2026}. The gaps are nuclear-grade heavy forgings, fuel fabrication \citep{Confindustria2025} and the staff of the safety authority, 79 people in 2026 \citep{ISIN2026}; Ansaldo itself indicates the need for a foreign partner as system integrator \citep{Ansaldo2026}.


\subsection{Political feasibility}
Law 169/2026 delegates the government to issue the implementing decrees within twelve months and does not authorise any plant \citep{Law2026}. The national plan foresees 8--16\,GW of nuclear in 2050, 11--22\% of demand \citep{PNIEC2024}. Compared in shares, our central case has <<RUN:ss0_PI_670_h3:N_GW|1>>\,GW, <<RUN:ss0_PI_670_h3:sh|0>>\% of demand, <<EXPR:('within','above')[min(1,int(V('RUN:ss0_PI_670_h3:N_GW')//16.0001))]>> the range of the plan (8--16\,GW, 11--22\%); on our curve (PI, six hours) the range of 8--16\,GW corresponds to fixN between <<THR6:PI:s:5:16|0>> and <<THR6:PI:s:5:8|0>>\,EUR/kW-yr, that is a construction cost of about <<EXPR:V('CPX:'+str(V('THR6:PI:s:5:16'))+':5:7')|-2>>--<<EXPR:V('CPX:'+str(V('THR6:PI:s:5:8'))+':5:7')|-2>>\,EUR/kW with 7 years of construction; the central cost (8{,}540\,EUR/kW, fixN 670) lies below this range, so building as at Dukovany makes more nuclear economic than the plan foresees. The 2011 referendum repealed the previous nuclear framework with 94\% of votes and 55\% turnout; some constitutional lawyers argue that reintroducing it by law is open to challenge \citep{Cabiddu2026}, while others consider the circumstances changed. Recent polls show a majority of about 55\% in favour of new-generation nuclear \citep{Euromedia2026}. The governing parties and Azione voted for the law; the Democratic Party, the Five Star Movement and AVS voted against \citep{Parties2026}. Regional laws against nuclear plants were struck down by the Constitutional Court in 2010 \citep{Consulta2010}. At EU level, nuclear is in the sustainable finance taxonomy, a challenge by Austria was dismissed in 2025 \citep{Taxonomy2025}, and state aid has been approved for comparable projects \citep{WNN2024}. A programme spanning at least three legislatures needs cross-party support: Germany, Belgium, Spain and Switzerland have changed nuclear policy at least once (Germany in 2010--2011 and 2022; Belgium in 2015, 2022 and 2025; Spain in 2011 and with the extension of Almaraz; Switzerland in 2011 and 2017, with a revision of the ban on new plants approved by Parliament in June 2026 and a referendum expected), although only Belgium has repealed its ban on new plants \citep{WNAcountries}.

\subsection{The role of solar PV}
At the real lifetime cost derived from the main FER~X auction, solar is a large part of every least-cost system, but not of the same kind. In the central case (PI, six hours) the system uses <<EXPR:V('GR:ss0_PI_670_h6:pv:GW')-V('UN:ss0_PI_670_h6:PVex_terra:GW')-V('UN:ss0_PI_670_h6:PVex_tetti:GW')|1>>\,GW of new solar, <<EXPR:V('UN:ss0_PI_670_h6:PV_terra_Sud-isole:GW')+V('UN:ss0_PI_670_h6:PV_terra-pipeline_Lazio/Sicilia/Piem./Puglia/Sard.:GW')+V('UN:ss0_PI_670_h6:PV_terra_Centro-Sud:GW')+V('UN:ss0_PI_670_h6:PV_terra_Nord:GW')|1>> ground-mounted in blocks that are all exhausted and <<UN:ss0_PI_670_h6:PV_agrivoltaico_Sud-isole:GW|1>> agrivoltaic, on top of the 24.2\,GW installed after 2016; without nuclear the PI system uses <<GR:ss0_PI_none_h6:pv:GW|0>>\,GW, of which <<EXPR:V('GR:ss0_PI_none_h6:pv:GW')-24.2|1>> new. Solar is therefore worth building next to nuclear, if it costs what the auctions pay and is not sized for the winter, which nuclear, wind and hydro cover. The plants that set the auction prices are large ground-mounted plants in the best sites. Residential rooftop systems costing 1{,}500--2{,}300\,EUR/kWp \citep{Rooftop2026} correspond to about 100--185\,EUR/MWh at 1{,}066--1{,}300 full-load hours, with 25-year life, 5\% cost of capital and 1.5\% annual O\&M (our calculation), two to four times the cost at which the model builds solar; their private profitability depends on tax credits and avoided network charges, which are transfers from other taxpayers and consumers. Before nuclear is available, every solar MWh displaces mostly gas: the auction price of 57\,EUR/MWh is about half the 2025 average wholesale price of 118\,EUR/MWh \citep{PUN2025}. Ground-mounted solar built now is therefore part of the least-cost path with or without nuclear.

\subsection{Limitations}\label{sec:limits}
The model has one node, no ramp constraints and demand response only through smart charging and interruptible load; it uses four historical weather years, the 2016--2019 load structure and a regime system for 2040, not a transition path. The adequacy reserve, the availability of nuclear (0.97 and 0.87), the shape of heat-pump and vehicle demand and the credits of hydro, pumped hydro and batteries are declared assumptions. Imports have the 2025 average price, exports earn nothing and no help comes from abroad in the critical weeks; the neighbouring systems are not simulated. The sensitivities on CSP, SMRs and gas with capture were computed under the earlier import rule. The import rule cuts both ways. Relative to the capacity-credit rule, net self-sufficiency raises the gap from <<GAP:val_PI_none:c3_PI_670_h6|2>> to <<GAP:ss0_PI_none_h6:ss0_PI_670_h6|2>>\,EUR/MWh (PI, six hours): not counting on other countries favours nuclear, and it is a stated policy choice, not a datum. Conversely, the model assumes that neighbours always absorb Italian exports up to the full interconnection capacity, even when they would also have surplus energy; this favours the system without nuclear, which exports <<UN:ss0_PI_none_h6:IMPORT:TWh|0>>\,TWh a year against <<UN:ss0_PI_670_h6:IMPORT:TWh|0>> in the central system. The nuclear share of demand includes the part of nuclear output that is exported. Hydropower is independent of the simulated weather because its monthly budgets are the same in the four years. The six-hour step underestimates the value of nuclear and absolute cost levels, and one-hour results cannot be computed from three-hour data; worst years exist at six hours only.

Several simplifications favour the system without nuclear: the single node, the absence of stability services, four weather years (with the worst years the gap grows) and technology costs for hydrogen and compressed air that are studies or targets. Others favour nuclear: availability of 0.92 and credit of 0.97 in the reserve, flat bioenergy and geothermal, an evening-heavy but otherwise inflexible demand, full credit for hydro and pumped hydro, and no ramp constraints. Existing solar at operating cost and the free wind block E0 favour renewables in both systems. The OB model is reduced, so OB and PI points differ by model as well as by cost case (without nuclear OB costs <<RUN:ss0_OB_none_h6:cost_new|2>> against <<RUN:ss0_PI_none_h6:cost_new|2>> for the full PI model, six hours); the SMR and wave tests switch those technologies back on, one at a time, on top of the OB base. Our own, unsourced assumptions are the 10\,GW CSP ceiling, the 20-year life and 20\% O\&M of 100-hour storage, the 3-year build of gas with CCS, the 1.287 grid scaling at 7\%, the DAC costs and the 17\% import mark-up. The central Dukovany value, 8{,}540\,EUR/kW, is the gross contract minus the contract fuel (itself disputed between sources) plus owner's costs taken from a US database (INL, 12.57\% of overnight cost; 15--20\% of the EPC cost in a European Commission synthesis \citep{DHaeseleer2013}), not from a Czech estimate; the EPC-only value (7{,}470) gives about <<N6:PI:s:5:604|0>>\,GW. The risk case, 9{,}680\,EUR/kW, is the median of four builds whose conversion to 2025 euros assumes uniform spending (Section~\ref{sec:methods}).

The unit size of the reactors (1.055 to 1.4\,GW) lies within the 10\% reserve margin (7.7\,GW at the three-hour net peak of 76.88\,GW), which covers the loss of five units. Common-mode failures across several units are not modelled.
